\section{Splitting a two-dimensional critical coefficient}
\label{cr11:sec:restriction}

Direct sum gives a restriction from a two-dimensional representation
of $\mathbb C[x,y,z]$ to two ordered point modules.
The vanishing-cycle coefficient supplies the restriction morphism.
An Euler class supplies its degree-zero normalization.
At coincident points, the coefficient contains an additional torsion
summand which the restriction annihilates.

All spaces have their complex analytic topology. Coefficients are rational.
For $n=1,2$, put
\[
 V_n=\operatorname{Mat}_n(\mathbb C)^3,\qquad
 f_n(A,B,C)=\operatorname{Tr}(A[B,C]),\qquad
 Z_n=\{[A,B]=[B,C]=[C,A]=0\}.
\]
The trace pairing identifies $Z_n$ with $\operatorname{Crit}(f_n)$.
Write $\mathfrak X_n=[V_n/\operatorname{GL}_n]$ and
$\mathfrak M_n=[Z_n/\operatorname{GL}_n]$.
Complexes on these stacks retain the full conjugation descent data.
An ordinary atlas pullback introduces no additional shift.
Our coefficient and grading are
\[
 \mathcal P_n=\left.\Phi_{f_n}\mathbb Q_{\mathfrak X_n}[2n^2]
                    \right|_{\mathfrak M_n},\qquad
 \Phi_fK=\operatorname{Cone}(i^*K\longrightarrow\psi_fK)[-1],
 \qquad H^j(K[r])=H^{j+r}(K).
\]
Here $i$ includes the zero fibre and $\psi_f$ denotes nearby cycles.
The fibre complex of $u:A\to B$ has differential
$d(a,b)=(da,u(a)-db)$.
In particular $\Phi_0=1$ and $\mathcal P_1=\mathbb Q[2]$.
Let
\[
 \mathcal H_n^k=\mathbb H^k(\mathfrak M_n,\mathcal P_n).
\]
Only ordinary hypercohomology and algebraic tensor products occur.
The function is a submersion away from $Z_n$, so its vanishing cycles
are supported on $Z_n$.
Extension by zero therefore identifies $\mathcal H_n$ with ambient
vanishing-cycle hypercohomology on $\mathfrak X_n$.

Let $T=(\mathbb C^*)^2\subset\operatorname{GL}_2$ be the diagonal torus.
The fixed subspace $D=V_2^T$ consists of diagonal triples and is
$\mathbb C^3\times\mathbb C^3$.
Thus $D\subset Z_2$ and
\[
 d:[D/T]=\mathfrak M_1\times\mathfrak M_1\longrightarrow\mathfrak M_2
\]
is the ordered direct-sum map.
Let $L_1,L_2$ be the two character lines of $T$, and put
\[
 R=H^*(BT;\mathbb Q)=\mathbb Q[x_1,x_2],\qquad
 x_i=c_1(L_i),\qquad |x_i|=2,\qquad \delta=x_2-x_1.
\]
These Chern classes are independent of the positions of the point modules.
The contraction of $D$ gives
$\mathcal H_1\otimes\mathcal H_1=R[4]$ as graded vector spaces.
A polynomial of ordinary degree $j$ has degree $j-4$ in this tensor product.
There is no additional torus acting on the three arrows.

The canonical first map of the vanishing-cycle triangle restricts to
\begin{equation}
 \epsilon_D:d^*\mathcal P_2\longrightarrow\mathbb Q_{[D/T]}[8].
 \label{cr11:eq:epsilon}
\end{equation}
This is a morphism, without an assertion that it is an isomorphism.
Equivalently, it is vanishing-cycle pullback for $f_2\circ d=0$.
The support construction of this pullback is
Davison\cite{cr11Davison}*{\S2.2, equations (19)--(23), pp.~10--11}.
The triangle itself defines \eqref{cr11:eq:epsilon} without a comparison theorem.
It induces
\begin{equation}
 \theta_{1,1}:\mathcal H_2^k\longrightarrow R^{k+8}
        =(\mathcal H_1\otimes\mathcal H_1)^{k+4}.
 \label{cr11:eq:theta}
\end{equation}
Its source contains the entire critical coefficient.

\section{The Euler factor of the extension correspondence}
\label{cr11:sec:euler}

Let $E\subset V_2\times\mathbb P^1$ classify triples preserving a line.
Write $p:[E/\operatorname{GL}_2]\to\mathfrak X_2$ for the proper map
forgetting that line.
The two diagonal blocks define
$q:[E/\operatorname{GL}_2]\to[D/T]$.
The potential $f_2$ is identically zero on $E$.
The section of $q$ with zero upper entries satisfies $p\circ s=d$.

In a frame with the submodule first, set
$W=\operatorname{Hom}(L_2,L_1)$.
The ambient extension stack has relative tangent complex
\begin{equation}
 [\,W\xrightarrow{\partial}W^{\oplus3}\,],\qquad
 \partial(t)=((a_2-a_1)t,(b_2-b_1)t,(c_2-c_1)t),
 \label{cr11:eq:extension}
\end{equation}
where $W$ has degree $-1$ and $W^{\oplus3}$ has degree zero.
The coordinates $(a_i,b_i,c_i)$ specify the two point modules.
Changing the sign of $\partial$ gives the same quotient action.
With the positive complex orientation, the Euler ratio is
\begin{equation}
 e_q=\frac{e_T(W^{\oplus3})}{e_T(W)}
       =\frac{(-\delta)^3}{-\delta}=\delta^2.
 \label{cr11:eq:euler}
\end{equation}
The calculation concerns the two-term complex on the ambient
three-loop representation stack.

The critical incidence is
$\mathcal F=[E/\operatorname{GL}_2]\times_{\mathfrak X_2}\mathfrak M_2$.
It retains the equations on upper entries.
If $v=(a_2-a_1,b_2-b_1,c_2-c_1)$ and $u$ is the upper-entry vector,
these equations are $v\wedge u=0$.
For $v\ne0$, they force $u\in\mathbb Cv$.
For $v=0$, they impose no condition on $u\in\mathbb C^3$.
Thus \eqref{cr11:eq:euler} comes from the ambient extension complex,
whose two vector bundles have constant ranks.

The map $q$ induces an isomorphism on ordinary cohomology with constant
coefficients: its upper-entry vector space and unipotent change-of-splitting
group are contractible.
Consequently $s^*=(q^*)^{-1}$ on this cohomology.
Naturality of vanishing-cycle pullback gives
\[
 \theta_{1,1}=s^*p^*=(q^*)^{-1}p^*,
\]
with the shifts in \eqref{cr11:eq:theta}.
Here $p^*$ is pullback from the ambient vanishing-cycle complex to
$\Phi_0\mathbb Q_E$, followed by cohomology.
Its construction retains the original critical support before applying
the coefficient morphism to $\mathbb Q_E$.
This identifies the map with the rank $(1,1)$ ambient restriction
in the localized Hall construction\cite{cr11JKL}*{\S3.2.3, Lemma~3.2.4,
pp.~26--27}.

\begin{definition}[Normalized rank $(1,1)$ restriction]
\label{cr11:def:delta}
The localized critical Hall restriction is the degree-zero map
\begin{equation}
 \Delta_{1,1}^{\mathrm{loc}}=\delta^{-2}\theta_{1,1}:
 \mathcal H_2\longrightarrow
       (\mathcal H_1\otimes\mathcal H_1)[\delta^{-1}].
 \label{cr11:eq:delta}
\end{equation}
Localization is algebraic localization of $R$ at powers of $\delta$.
It is not a completion. No coefficient field extension occurs.
\end{definition}

The inverse Euler class has degree $-4$, which cancels the degree
$+4$ of \eqref{cr11:eq:theta}.
The scalar calculation below shows that this particular map has
polynomial image.

\section{The Hall identity for two lines}
\label{cr11:sec:hall}

The embedding $E\subset V_2\times\mathbb P^1$ has complex codimension three.
Its complex Thom map, followed by integration over $\mathbb P^1$, is
\[
 \gamma:R\bar p_*\mathbb Q_E\longrightarrow\mathbb Q_{V_2}[4],
 \qquad \bar p:E\longrightarrow V_2.
\]
All maps descend under $\operatorname{GL}_2$.
Applying $\Phi_{f_2}$, using proper exchange and $f_2|_E=0$, then
restricting to $Z_2$, gives
\[
 \mu_{1,1}:Rb_*\mathbb Q_{\mathcal F}[4]\longrightarrow\mathcal P_2,
 \qquad b:\mathcal F\longrightarrow\mathfrak M_2.
\]
Pullback from the two diagonal blocks and this morphism define
$m_{1,1}:\mathcal H_1\otimes\mathcal H_1\to\mathcal H_2$.
The order is submodule followed by quotient.

\begin{theorem}[Restriction of a rank $(1,1)$ product]
\label{cr11:thm:hall}
For $a,b\in\mathbb Q[x]=H^*(\mathbb C^3\times B\mathbb C^*;\mathbb Q)$,
regarded with shift $[2]$, one has
\begin{align}
 \theta_{1,1}m_{1,1}(a,b)
   &=\delta^2\bigl(a(x_1)b(x_2)+a(x_2)b(x_1)\bigr),
       \label{cr11:eq:raw-hall}\\
 \Delta_{1,1}^{\mathrm{loc}}m_{1,1}(a,b)
   &=a\otimes b+b\otimes a.
       \label{cr11:eq:hall}
\end{align}
In particular, if $h_2=m_{1,1}(1,1)$ and $\sigma_2=h_2/2$, then
\begin{equation}
 \theta_{1,1}(\sigma_2)=\delta^2,\qquad
 \Delta_{1,1}^{\mathrm{loc}}(\sigma_2)=1\otimes1.
 \label{cr11:eq:sigma}
\end{equation}
\end{theorem}

\begin{proof}
Naturality of the first map $\Phi_fK\to i^*K$ in the coefficient $K$
commutes the triangle with $\gamma$.
Proper exchange commutes this square with the flag direct image.
Since the potential on $E$ is zero, its triangle map is the identity.
Thus the composite on the left of \eqref{cr11:eq:raw-hall} is the
ordinary complex-oriented incidence Gysin map, restricted to $D$.
This comparison is made before taking cohomology.

The two $T$-fixed flags over $D$ are the coordinate lines $L_1,L_2$.
At $L_1$, the three forbidden lower entries have weight $\delta$.
The flag tangent has weight $\delta$.
At $L_2$, both weights change sign.
Their contributions after localization are therefore
\[
 \frac{\delta^3}{\delta}a(x_1)b(x_2),\qquad
 \frac{(-\delta)^3}{-\delta}a(x_2)b(x_1).
\]
Both coefficients are $+\delta^2$.
The complex orientation determines these signs.
The usual equivariant self-intersection and integration formulas apply
to the smooth ambient embedding and the projective flag fibre.
Their compatibility with critical pullback is also the one in
Davison\cite{cr11Davison}*{\S2.7, Proposition~2.16, pp.~21--22}.
The expression is polynomial, and $R\to R[\delta^{-1}]$ is injective.
Hence the equality already holds in $R$.
Division gives \eqref{cr11:eq:hall} and \eqref{cr11:eq:sigma}.
The rank-one groups have even degree, so the exchange introduces
no Koszul sign.
\end{proof}

There is also a calculation without fixed-point fractions.
Let $S$ be the tautological line on $\mathbb P(L_1\oplus L_2)$,
write $\xi=c_1(S^*)$ and put $c_1=x_1+x_2$, $c_2=x_1x_2$.
Then
\[
 \xi^2+c_1\xi+c_2=0,\qquad
 \int_{\mathbb P^1}\xi=1,\qquad
 c_1(\operatorname{Hom}(S,(L_1\oplus L_2)/S))=c_1+2\xi.
\]
The coefficient of $\xi$ in $(c_1+2\xi)^3$ modulo this relation is
$2(c_1^2-4c_2)=2\delta^2$.
This gives $\theta_{1,1}(h_2)=2\delta^2$ with the same complex orientation.

\section{The full scalar coefficient and its restriction}
\label{cr11:sec:scalar}

Let $s=(aI_2,bI_2,cI_2)$ be a scalar triple.
Scalar translation preserves the potential, so take $s=0$.
Let $P_s$ be the ordinary atlas stalk of $\mathcal P_2$ and define
\[
 K_s=R\Gamma(BT,P_s).
\]
This notation retains the equivariant derived stalk, including its
extension data. It does not mean the Borel cohomology of a direct
sum of ordinary stalk cohomology groups.
Restriction of \eqref{cr11:eq:epsilon} gives
\[
 \theta_s:H^k(K_s)\longrightarrow R^{k+8}.
\]

\begin{theorem}[Scalar Euler image and torsion]
\label{cr11:thm:scalar}
The equivariant scalar coefficient has a decomposition of graded
$R$-modules
\begin{equation}
 H^*(K_s)=R\,v\ \oplus\ (R/(\delta^2))\,w,
 \qquad |v|=-4,\qquad |w|=-2,
 \label{cr11:eq:scalar-module}
\end{equation}
where $v$ is the scalar restriction of $\sigma_2$.
The map is
\begin{equation}
 \theta_s(v)=\delta^2,\qquad \theta_s(w)=0.
 \label{cr11:eq:scalar-map}
\end{equation}
The choice of $w$ depends on a nonzero degree-five Milnor class.
The free generator $v$ has the Hall normalization above.
The ordinary atlas coefficient has one-dimensional groups in degrees
$-4,-2,1$, and no other groups.
\end{theorem}

\begin{proof}
Write the traceless parts in the ordered basis
$H=\operatorname{diag}(1,-1)$, $U=e_{12}$, $L=e_{21}$.
If $M$ has the three coordinate columns, direct multiplication gives
\begin{equation}
 f_2=2\det M.
 \label{cr11:eq:det}
\end{equation}
The three scalar coordinates are free.
The compact torus $K=(S^1)^2$ acts on $M$ by left multiplication
through $\operatorname{diag}(1,t_1/t_2,t_2/t_1)$.

Choose a $K$-invariant Hermitian norm.
The support description of normalized vanishing cycles uses
$P=\{\operatorname{Re}f_2\leq0\}$ and the fibre of restriction from
a small ball to its complement of $P$.
This description is Massey\cite{schMasseyST}*{\S1, p.~2}.
The radial map $x\mapsto x/(1+\|x\|/r)$ gives equivariant homotopy
equivalences from the global pair to the pair in the radius-$r$ ball.
It preserves $P$, since the potential is homogeneous of degree three.
On $H_+=\{z:\operatorname{Re}z>0\}$ use the cube root positive on
the positive real axis. The equivariant homeomorphism
\[
 V_2\setminus P\simeq f_2^{-1}(1)\times H_+,\qquad
 x\longmapsto(f_2(x)^{-1/3}x,f_2(x))
\]
identifies the two restriction maps.
Thus there is an actual fibre sequence
\begin{equation}
 K_s\longrightarrow R\Gamma(BK;\mathbb Q)[8]
   \longrightarrow R\Gamma(f_2^{-1}(1)\mathbin{/\mkern-6mu/}K;\mathbb Q)[8],
 \label{cr11:eq:borel-triangle}
\end{equation}
where $X\mathbin{/\mkern-6mu/}K=EK\times_KX$ denotes the Borel construction.
The first arrow is $\theta_s$ on cohomology.

After constant scaling, the nonzero determinant fibre is
$\operatorname{SL}_3(\mathbb C)$.
Polar decomposition $M=UP$ gives the deformation $UP\mapsto UP^t$
onto $\operatorname{SU}(3)$.
It is equivariant for the displayed unitary left action.
The scalar-coordinate factor is equivariantly contractible.
Split $K=K_c\times K_e$ by
$(z,t)\mapsto(zt,z)$.
The central factor $K_c$ acts trivially, while
\[
 K_e\longrightarrow\operatorname{SU}(3),\qquad
 t\longmapsto\operatorname{diag}(1,t,t^{-1})
\]
acts freely by left multiplication.
Consequently the last Borel space in \eqref{cr11:eq:borel-triangle}
has the homotopy type
\[
 BK_c\times Y,\qquad Y=K_e\backslash\operatorname{SU}(3).
\]

The map taking the first row of a unitary matrix gives a bundle
$\operatorname{SU}(2)\to\operatorname{SU}(3)\to S^5$.
Its structure group acts by left multiplication in the lower
$2\times2$ block, which contains $K_e$.
Taking the left $K_e$ quotient gives
\begin{equation}
 S^2=K_e\backslash\operatorname{SU}(2)\longrightarrow Y\longrightarrow S^5.
 \label{cr11:eq:sphere-bundle}
\end{equation}
Local orthonormal completion of a first row supplies local trivializations.
Changes of completion induce right translations on
$K_e\backslash\operatorname{SU}(2)$, preserving orientation.
The rational Serre spectral sequence of \eqref{cr11:eq:sphere-bundle}
has terms only in bidegrees $(0,0),(0,2),(5,0),(5,2)$.
Every differential vanishes for degree reasons.

Put $u=c_1(t)=x_1-x_2=-\delta$.
The principal circle bundle
$\operatorname{SU}(3)\to Y$ restricts on the displayed $S^2$ fibre
to $S^3\to S^2$.
In $\operatorname{SU}(2)$ coordinates
$\left(\begin{smallmatrix}a&b\\-\bar b&\bar a\end{smallmatrix}\right)$,
left multiplication by $\operatorname{diag}(t,t^{-1})$ sends
$(a,b)$ to $(ta,tb)$.
Thus its quotient is the Hopf map $S^3\to\mathbb P^1$,
whose transition function on the equator has degree one up to orientation.
Thus $u$ restricts to a generator of $H^2(S^2;\mathbb Q)$, up to
its orientation sign.
Let $\beta\in H^5(Y;\mathbb Q)$ be the pullback of a sphere generator.
Fibre integration shows $u\beta\ne0$.
There is no degree-four cohomology, and odd-degree squares vanish over
$\mathbb Q$. Therefore
\[
 H^*(Y;\mathbb Q)=\mathbb Q[u]/(u^2)\otimes\Lambda(\beta),
 \qquad |\beta|=5.
\]
The homomorphism induced by $Y\to BK_e$ sends $u$ to this class.
After adjoining the central variable this gives
\begin{equation}
 H^*(f_2^{-1}(1)\mathbin{/\mkern-6mu/}K;\mathbb Q)
       =R/(\delta^2)\otimes\Lambda(\beta).
 \label{cr11:eq:borel-fibre}
\end{equation}
The map from $R$ is the quotient into the even summand.

The long exact sequence of \eqref{cr11:eq:borel-triangle} now gives
\[
 0\longrightarrow (R/(\delta^2))\,w
 \longrightarrow H^*(K_s)
 \xrightarrow{\theta_s}\delta^2R[8]
 \longrightarrow0,
 \qquad |w|=5+1-8=-2.
\]
The quotient is a free graded $R$-module with generator in degree $-4$.
Equation~\eqref{cr11:eq:sigma}, restricted to the scalar diagonal,
supplies a section $v$ whose image is $\delta^2$.
This proves \eqref{cr11:eq:scalar-module} and \eqref{cr11:eq:scalar-map}.

For the ordinary atlas stalk, the bundle
$S^3\to\operatorname{SU}(3)\to S^5$ gives
$H^*(\operatorname{SU}(3);\mathbb Q)=\Lambda(\alpha_3,\alpha_5)$.
Its spectral sequence has no possible differential.
The product of the two generators evaluates nontrivially on the total
orientation by fibre integration.
The reduced cohomology degrees are $3,5,8$.
The normalized fibre and shift $[8]$ send these degrees to $-4,-2,1$.
In degree $-4$, the equivariant-to-ordinary edge map is an isomorphism:
there is no lower stalk cohomology which could contribute or receive
a differential in that total degree.
Hence $v$ restricts to a nonzero ordinary lowest class.
\end{proof}

The scalar Borel spectral sequence gives a direct description of the
extension responsible for the torsion.
Normalize $\alpha_3$ so that
\[
 d_4\alpha_3=\delta^2,\qquad
 d_4(\alpha_3\alpha_5)=\delta^2\alpha_5.
\]
The first differential is nonzero because
\eqref{cr11:eq:borel-fibre} has no degree-four class in the effective-circle
variable. Its sign fixes the chosen transgression normalization.
The second equality is the product rule.
Thus the ordinary odd scalar class participates in the equivariant
coefficient, although $H^*(K_s)$ is concentrated in even degrees.
Taking the tensor product of \eqref{cr11:eq:scalar-module} with
$R/(x_1,x_2)$ does not compute the ordinary atlas stalk.
The original fibre complex in \eqref{cr11:eq:borel-triangle} retains
that distinction.

Potential monodromy also retains its ambient meaning.
The maps $(A,B,C)\mapsto(A,e^{i\vartheta}B,C)$ preserve centered balls
and multiply the potential by $e^{i\vartheta}$.
At $\vartheta=2\pi$ the map is the identity.
These maps commute with conjugation and the constant-to-nearby-cycle map.
They give identity monodromy on the scalar equivariant fibre complex.
The triangle morphisms commute with this monodromy.

\section{Polynomial image and specialization}
\label{cr11:sec:polynomial}

\begin{corollary}[Cancellation in rank two]
\label{cr11:cor:polynomial}
The map \eqref{cr11:eq:delta} has image in
$\mathcal H_1\otimes\mathcal H_1$.
It therefore defines a unique degree-zero map
\[
 \Delta_{1,1}:\mathcal H_2\longrightarrow\mathcal H_1\otimes\mathcal H_1,
 \qquad \theta_{1,1}=\delta^2\Delta_{1,1}.
\]
It satisfies \eqref{cr11:eq:hall} and \eqref{cr11:eq:sigma}.
Its scalar counterpart sends $v$ to $1\otimes1$ and annihilates
$(R/(\delta^2))w$.
\end{corollary}

\begin{proof}
Include one scalar point in $D$ and retain its $T$ stabilizer.
Restriction $H^*([D/T],\mathbb Q)\to H^*(BT,\mathbb Q)$ is the
identity on $R$, since $D$ is equivariantly contractible.
Naturality of \eqref{cr11:eq:epsilon} identifies the scalar restriction
of $\theta_{1,1}(c)$ with $\theta_s$ of the restriction of $c$.
Theorem~\ref{cr11:thm:scalar} puts this image in $\delta^2R$.
Therefore $\theta_{1,1}(c)$ is divisible by $\delta^2$.
Multiplication by $\delta^2$ is injective on $R$.
This proves existence and uniqueness of the polynomial map.
The scalar statements follow from \eqref{cr11:eq:scalar-map}.
\end{proof}

Collision of the two point positions means restriction along the scalar
diagonal $\mathbb C^3\to D$.
This leaves $x_1,x_2$ and $\delta$ as stabilizer classes.
The natural coefficient square for this restriction is the square
used in the corollary.
It pulls back the ambient vanishing cycles rather than replacing them
by vanishing cycles of a zero function on the scalar locus.

Forgetting stabilizers is a different operation.
It sends $x_1,x_2$ to zero.
There is no ring homomorphism $R[\delta^{-1}]\to\mathbb Q$ extending
that specialization.
After the proved cancellation, evaluation of the polynomial output is
well defined. It sends the scalar Hall generator $v$ to $1$.
The raw ordinary stalk morphism
$P_s\to\mathbb Q[8]$ is zero: its target degree is $-8$, whereas
the source degrees are $-4,-2,1$.
Consequently evaluation of the cancelled cohomological map and
ordinary pullback of the raw coefficient morphism are distinct operations.

Interchange of $L_1,L_2$ sends $\delta$ to $-\delta$ and preserves
$\delta^2$. The two direct-sum maps differ by conjugation in
$\operatorname{GL}_2$.
Naturality therefore makes $\Delta_{1,1}$ symmetric under this interchange.
Its scalar generator is the complete-flag class $h_2/2$.
This is also the input normalization for products involving two
rank-two constituents.
For a flag of type $(2,2,1)$, the restriction constructed here can
be applied to either rank-two input and sends its lowest class to
the ordered pair of rank-one classes.
A coproduct identity for the rank-five output would additionally
require restriction maps with rank-five source and a comparison of
the corresponding flag correspondences.
The map above establishes the rank $(1,1)$ identity without that
additional assertion.
